\documentclass[../../../maim.tex]{subfiles}

\begin{document}


\begin{center}
\includegraphics[width = \textwidth]{lattice-0-flux-state.jpeg}
\end{center}

This system is a simple triangular Bravais lattice with primitve translations

\begin{equation}
	\vec{a}_1 = 
	\begin{pmatrix}
		1 \\
		0
	\end{pmatrix}, \quad 
	\vec{a}_2 = 
	\begin{pmatrix}
		1/2 \\
		\sqrt{3}/2
	\end{pmatrix}.
\end{equation}

\subsubsection{Pedestrian method}

The tight-binding Hamiltonian is:

\begin{align}
	H &=  \sum_{i, \delta, \sigma} \frac{t}{2} \left(f_{i\sigma}^\dagger f_{i+\delta \, \sigma} + \text{h.c.}\right) + \mu \sum_{i}\left(f_{i\uparrow}^\dagger f_{i\uparrow} + f_{i\downarrow}^\dagger f_{i\downarrow} - 1\right) + \text{const.} \notag \\
	&= \sum_{i, \delta, \sigma} \frac{t}{2} \left(f_{i\sigma}^\dagger f_{i+\delta \, \sigma} + \text{h.c.}\right) + \mu \sum_{i, \sigma}f_{i\sigma}^\dagger f_{i\sigma} - \mu N + \text{const.}.
\end{align}

$i$ runs over sites sites and $\delta$ over nearest neighbours of a generic site. I include the absolute shift for the sake of consistency. The factor 1/2 is also included for consistency. The constant is a product of mean field theory and reads

\begin{equation}
	\frac12\sum_{i,\delta} 2\tilde{J}\left(1-6\tilde{\kappa}\abs{\langle\chi_{i,i+\delta}}^2\rangle\right)\abs{\langle\chi_{i,i+\delta}\rangle}^2 = 6N\left( \tilde{J}\left(1-6\tilde{\kappa}\abs{\langle\chi_{i,i+\delta}\rangle}^2\rangle\right)\abs{\langle\chi_{i,i+\delta}\rangle}^2\right)
\end{equation}

The second step can be understood in at least two ways; a) each site has 6 neighbours, or b) there are  bond per primitive unit cell 1 of which is completely inside the unit cell and contributes a 2 and the other 4 shared between 2 unit cells thereby contributing a 1. All bonds are equivalent of course. The afformentioned six nearest neighbours lie at the (relative) positions

\begin{equation}
\vec{\delta}_1=
\begin{pmatrix}
1 \\ 0
\end{pmatrix} \; 
\vec{\delta_2}=
\begin{pmatrix}
\frac{1}{2} \\ \frac{\sqrt{3}}{2}
\end{pmatrix} \;
\vec{\delta_3}=
\begin{pmatrix}
\frac{1}{2} \\ -\frac{\sqrt{3}}{2}
\end{pmatrix}
\notag 
\end{equation}

and $-\vec{\delta}_1, \; -\vec{\delta}_2, \; -\vec{\delta}_3$. I will start with the $0-0$ state. To advance, do a Fourier transform:

\begin{equation}
f_{i \sigma} = \frac{1}{\sqrt{N}}\sum_k e^{i\vec{k}\vec{R}_i} c_{k \sigma}.
\end{equation}

\begin{align}
	\Rightarrow H &= \sum_{q,\sigma} \frac{t}{2} \left(c_{q\sigma}^\dagger c_{q\sigma} \left(\sum_\delta e^{i\vec{q}\vec{\delta}} + \mu \right) + \text{h.c.}\right) - \mu N \notag \\
	&= \sum_{q,\sigma}\left(t\sum_{\delta} \cos(\vec{q}\vec{\delta}) + \mu\right) c_{q\sigma}^\dagger c_{q\sigma} - \mu N
\end{align}

Using first the symmetrie of the cosine and then the sum identity for the cosine the dispersion reads

\begin{align}
	\varepsilon(\vec{q}) &= 2t\left[\cos(q_x) + \cos(\frac12(q_x + \sqrt{3}q_y)) + \cos(\frac12(q_x - \sqrt{3}q_y))\right] + \mu \notag \\
&= 2t \left[\cos(q_x) + 2 \cos(\frac{q_x}{2})\cos(\frac{\sqrt{3}q_y}{2})\right] + \mu
\end{align}

features:

\begin{itemize}
	\item $\varepsilon(\vec{q}) \; \in \; [-3t, 6t] + \mu$ $\Rightarrow$ $\mu \; \in \; [-6t, 3t]$
	\item symmetries $q_x \rightarrow -q_x$ and $q_y \rightarrow -q_y$ (+ should have some rotational symmetrie and (thereby associated) inflections along other axes) \textcolor{red}{could check this}
	\item periodicity: $4\pi$ in $q_x$, $(4\pi)/\sqrt{3}$ in $q_y$
\end{itemize}

The mean field parameters are

\begin{equation}
	\langle \chi_{ij} \rangle  = \frac{1}{2} \langle f_{i \uparrow}^\dagger f_{j \uparrow} + f_{i \downarrow}^\dagger f_{j \downarrow} \rangle = \langle f_{i \uparrow}^\dagger f_{j \uparrow} \rangle
\end{equation}

Simplify the expression by Fourier transform:

\begin{align}
	\langle \chi_{ij} \rangle &= \frac{1}{N} \sum_{q,p} e^{i(\vec{q}\vec{R}_j - \vec{k}\vec{R}_i)} \langle c_p^\dagger c_q \rangle \notag \\
	&= \frac{1}{N}\sum e^{i\vec{q}(\vec{R}_j-\vec{R}_i)}f(\varepsilon(\vec{q}))
\end{align}

Take the thermodynamic limit $\sum_k \rightarrow (N\Omega)/(2\pi)^2\int dk$ ($\Omega$ is the unit cell volume).

\begin{align}
	\Rightarrow \langle \chi_{ij} \rangle = \Omega\int\limits_{\text{1.BZ}} \frac{d^2q}{(2\pi)^2} e^{i\vec{q}\left(\vec{R}_j - \vec{R}_i\right)}f(\varepsilon(\vec{q})) 
\end{align}

\textcolor{red}{It looks like the model allows for different $\chi_{ij}$ in each direction $\vec{R}_j - \vec{R}_i$ (there are 3 of those)?} \\

From this expression, recover the on-site particle number by setting $i=j$ and accounting for the factor $1/2$ in the definition of $\chi_{ij}$:

\begin{equation}
\langle N_i \rangle = 2\Omega\int \frac{d^2q}{(2\pi)^2} f(\varepsilon(\vec{q})).
\end{equation}

By construction the particle number is uniform. 

Some elaborations on the integral domain follow. The primite translations of the reciprocal lattice are
\begin{equation}
	\vec{b}_1=2\pi
	\begin{pmatrix}
		0 \\
		2/\sqrt{3}
	\end{pmatrix}, \quad
	\vec{b}_2 = 2\pi
	\begin{pmatrix}
		1 \\
		-1/\sqrt{3}
	\end{pmatrix}
\end{equation}

I do a basis transformation onto thos axes defined by 

\begin{gather}
	\begin{pmatrix}
		k_x \\
		k_y
	\end{pmatrix} =
	\begin{pmatrix}
		1/2 & \sqrt{3}/2 \\
		1 & 0 
	\end{pmatrix}
	\begin{pmatrix}
		q_x \\
		q_y
	\end{pmatrix} \\
	\Leftrightarrow 
	\begin{pmatrix}
		q_x \\
		q_y
	\end{pmatrix}=
	\begin{pmatrix}
		0 & 1 \\
		2/\sqrt{3} & -1/\sqrt{3} 
	\end{pmatrix}
	\begin{pmatrix}
		k_x \\
		k_y
	\end{pmatrix}\\
	\Rightarrow \text{det$J$} = -2/\sqrt{3} \\
	\Rightarrow \varepsilon(\vec{k}) = 2t\left(\cos(k_x) + \cos(k_y) + \cos(k_x -k_y)\right) + \mu
\end{gather}

The first Brillouin zone is $\{k_x, k_y\} \in [-\pi,\pi]^2$. The unit cell volume is $\Omega = \sqrt{3}/2$. Consistency check: $|\text{det}($J$)|$ = $\Omega^{-1}$, since the substitution simply transforms the unti cell to a square of unit volume. So the integrals that have to be evaluated are of the form

\begin{gather}
	\langle N_i \rangle = 2\int\limits_{[-\pi,\pi]^2} \frac{d^2k}{(2\pi)^2} f(\varepsilon(\vec{k}), \mu), \\
	\langle \chi_{ij} \rangle = \int\limits_{[-\pi,\pi]^2} \frac{d^2k}{(2\pi)^2} \exp[i\left(k_y, (1/\sqrt{3})(2k_x-k_y)\right)\cdot(\vec{R}_j-\vec{R}_i)] f(\varepsilon(\vec{k}), \mu)
\end{gather}

\textcolor{red}{Check that under a particle hole transformation chi changes sign, sign of t is doing just that and this brings the $pi$-flux phase.}

The condition is $\langle N_i \rangle \overset{!}{=} 1 \; \forall i$. By setting setting $t = -1$, the chemical potential is numerically determined in units of $t$ using mathematica. I choose $t=-1$ because I anticipate $t$ to be negative for a positive $\chi$. (This may not be true depending on $\kappa$ and $\chi$, but as long as it is self-consistent evrything is fine.) For the numerical integration I used a Fermi-function with a small temperature ($\beta = 200t$) instead of a step-function, since the ladder is not analytical and causes problems for the numerical algorithm. The result is $\mu = -0.834737 \; t$. The bond-parameters are calulcated the same way, the result is $\langle \chi_{ij} \rangle = 0.164713$. \\

Now, it is possible to express $\mu$ only in units of $J$. Use the formula

\begin{equation}
	t_{ij} = -J \langle\chi_{ij}\rangle^* (1 + \kappa/2 - 4 \kappa \left|\langle\chi_{ij}\rangle \right|^2).
\end{equation}

Set $\kappa = 10$:

\begin{gather}
	t_{ij} = -0.8095 \; J \\
	\mu = -0.67574 \; J \\
	\langle \chi_{ij} \rangle = 0.164713
\end{gather}

The energy density (that is also the free energy desnity at $T=0$) is  

\begin{equation}
	\frac{\langle H \rangle}{N} = 2\int\limits_{[-\pi,\pi]^2}\frac{d^2k}{(2\pi)^2} \varepsilon(\vec{k}) f(\varepsilon(\vec{k})) - \mu + 6 \tilde{J}\left(1-6\tilde{\kappa}\abs{\langle\chi_{i,i+\delta}\rangle}^2\rangle\right)\abs{\langle\chi_{i,i+\delta}\rangle}^2
\ = -0.888304 \, J
\end{equation}

Observtaion; the chemical potential drops out of the Free energy. \\ 

One can also directly calculate results for the state with all fluxes $\pi$ by performing a particle-hole-transformation:

\begin{gather}
	t \rightarrow -t \notag \\
	\mu \rightarrow -\mu \\
	\Rightarrow \varepsilon(\vec{k}) \rightarrow -\varepsilon(\vec{k}) \\
	\Rightarrow f(\varepsilon(\vec{k})) \rightarrow 1-f(\varepsilon(\vec{k})
\end{gather}

It can be shown that

\begin{equation}
\chi \rightarrow -\chi.
\end{equation}

The result is:

\begin{equation}
	\frac{\langle H \rangle_\pi}{N} = -0.888302 \, J
\end{equation}

\subsubsection{Pedro's method}

The longer explenation for this method is in the next section. Here I will only say that the free energy is the same except for $\gamma$, which gets replaced by a factor $\phi$ (as mentioned above the chemical potentioal does not contribute to the free energy):

\end{document}
